\chapter*{Conclusion}
\addcontentsline{toc}{chapter}{Conclusion}
\label{chap:conclusion}

This thesis set out to improve a genetic-programming hyper-heuristic (GPHH)
baseline for the multi-mode resource-constrained project scheduling
problem, and pursued that goal through a two-stage methodology: a broad,
cheap exploratory screen of \changed{candidate algorithmic modifications}
(Section~\ref{sec:parameter-sensitivity}), followed by a fully powered,
comprehensive evaluation on MMLIB50 and, as a full-scale scalability check,
MMLIB100 (Section~\ref{sec:operator-evaluation}). \changed{Of the screened variants,
only epsilon-lexicase selection produced a positive signal that survived
dedicated, escalating follow-up experiments; the others, including all
three diagnosis-driven, graft-based DMGE-family variation schemes, did not
outperform the baseline and were set aside, verdicts an
independent-seed replication subsequently confirmed.} Against an
already-tuned, published baseline this is an unsurprising base rate, and
the screen's value lay in establishing it cheaply before the comprehensive
budget was committed, leaving lexicase selection to be carried forward on
the strength of its dedicated experiments, alongside two independently
motivated extensions: critical-path local search and non-renewable (NR)
resource terminals.

\changed{The comprehensive stage produced three distinct results.} NR
terminals deliver a large, statistically significant improvement in both
schedule quality and feasibility across the entire MMLIB50 grid, and the
effect replicates on MMLIB100 in \changed{four of the six combinations of
schedule generation scheme and decision strategy}, losing significance
only under parallel/activity-first and
parallel/mode-first. This is the thesis's strongest and most consistently
replicated result: giving the GP explicit visibility into a resource
constraint it previously could not see converts a substantial share of
otherwise-infeasible schedules into feasible ones while also improving
quality. \changedB{A final re-evaluation of every evolved rule on the
complete 108-class held-out test set (Section~\ref{subsec:full108})
sharpens this verdict in one respect: the schedule-quality advantage holds
in all twelve benchmark cells at full test scale, but the feasibility gain
proves largely confined to the instance classes seen during training, so
the quality improvement is the part of the result that generalises.}
Lexicase selection reliably improves training-set fitness on both
benchmark sizes, but the advantage does not transfer reliably to held-out
test fitness, and under serial SGS it can reverse into a significant loss,
more pronouncedly so on MMLIB100; the most plausible reading is
selection-induced overfitting, and this thesis's contribution is to show
that the trade-off is not just present but SGS-dependent and strengthens
under scaling. \changedC{Critical-path local search does improve the
baseline at MMLIB50's instance size: it produces a significant win under
parallel SGS and the best average test-fitness rank of the four jointly
ranked conditions, so for problems at that scale it is a worthwhile
extension. Its benefit is, however, specific to that scale: on MMLIB100 the
significant improvement disappears and its rank falls from best to worst.}
\changed{This comparison shows directly that a fully powered
result obtained on one benchmark cannot be assumed to hold at another
scale until it has been checked there.} Combining lexicase and local search naively, finally,
underperforms both components individually and is not recommended by this
evidence.

Read against the research questions: a staged,
exploratory-then-comprehensive methodology did reliably narrow a large
field of candidates to a small set worth rigorous evaluation, with the
caveat that a cheap screen speaks to its own precision, not to whether it
missed a genuine effect among the ideas it ruled out. Epsilon-lexicase
selection does not, on this combined evidence, improve generalised schedule
quality across the tested grid, despite a real and scale-persistent
training-fitness advantage. \changedC{Critical-path local search improves
generalised schedule quality on MMLIB50-sized instances, but the advantage
does not carry to MMLIB100.} Exposing non-renewable resource state
to the GP terminal set does improve schedule quality relative to a matched
baseline, robustly across scale for four of six strategy/SGS combinations.

This evidence carries the usual qualifications of a controlled
computational study. Internal validity is protected by every condition
sharing one code base and differing only in the mechanism under test,
though the stratified ten-class subsample ties the absolute numbers to
that particular sample. External validity is where the evidence is most
direct: two instance scales sufficed to show that scale-independence is
not a safe default for any of the three mechanisms, but not to establish a
trend, and generalisation beyond the MMLIB family is not established here.
\phantomsection\label{subsec:threats-construct}
Construct validity rests on ARD\% against an unconstrained CPM lower
bound, which is systematically looser on the NR-preserving instances than
on the renewable-only ones, so ARD\% magnitudes are never compared across
that instance-set boundary.
\phantomsection\label{subsec:threats-stat}
Statistical conclusion validity is qualified by the absence of a
multiple-comparisons correction across the thirty-six pairwise Wilcoxon
tests, under which one or two isolated significant cells are expected by
chance alone, and by the joint-ranking Friedman test's low block count,
which limits its power and makes its MMLIB100 non-significance weak
evidence of no difference.

Three broader implications follow. For scheduling research, a
resource-visibility gap in a hyper-heuristic's terminal set is worth
checking for and closing before reaching for elaborate search-side
modifications: the largest effect here came from giving the GP information
it lacked, not from changing how it searched. For genetic programming more
broadly, the SGS-dependent lexicase gap is a concrete data point in the
ongoing conversation about selection pressure and generalisation, worth
checking whenever lexicase is applied to a multi-instance training regime.
For experimental methodology, the two-stage design with an explicit
scalability check is offered as a template for GPHH studies facing a large
candidate space under a bounded compute budget.

\phantomsection\label{sec:future-work}
Several follow-ups are motivated directly by what this evidence left open;
both benchmarks are already evaluated at full scale, so all of them are
about explaining or extending what the evidence shows. The most
consequential open question is why local search's MMLIB50 benefit
vanished at MMLIB100, best answered by instrumenting the hill-climb's
accepted-move statistics on both benchmarks side by side. The two
non-replicating NR-terminal cells invite a direct test of the split-tree
coordination account proposed in Section~\ref{subsec:nr-eval}, for
instance via a decision-strategy-specific NR terminal. Evaluating on
MMLIB+, whose tighter non-renewable budgets are exactly the setting NR
terminals target, would extend the preliminary check already made there.
Local search could be made SGS-adaptive once instrumentation shows where
its budget is productive. The DMGE family's two candidate failure
explanations, category coarseness and graft disruptiveness, remain to be
distinguished, cheaply screenable with the same pilot methodology.
Lexicase's generalisation gap deserves direct study via per-seed
training/validation trajectories. And the hybrid's underperformance rules out
one naive combination, not the idea of combining selection-side and
search-side mechanisms more selectively.

Many ideas were investigated over the course of this thesis, and most did
not improve on the baseline. Systematic, staged experimentation, rather
than intuition or a single exhaustive sweep, is what separated the
directions genuinely worth pursuing from a much larger field of
plausible-looking alternatives, and what allowed one of them, non-renewable
resource terminals, to be confirmed as the thesis's strongest and most
consistent result. \changed{The final contribution is a carefully evaluated
and reproducible set of findings: one extension that robustly improves the
baseline across benchmark scales, and a documented account of which other
directions do not.}
